% Generated by tools/regs.py from data/registers/interrupts.toml. Do not edit.
\begin{description}
    \item[INT\_PENDING\_0] (\verb+$FE20+, read/write, reset \verb+$00+) Group 0 pending flags. A bit reads 1 when its source has raised an interrupt. Write 1 to a bit to clear it; writing 0 leaves it alone.
    \begin{description}
        \item[INT\_CARTRIDGE] (bit 7) Interrupt line on the cartridge connector.
        \item[INT\_DMA] (bit 6) DMA transfer finished (when enabled in the DMA control register).
        \item[INT\_TIMER\_1] (bit 5) Timer 1 reached its target value.
        \item[INT\_TIMER\_0] (bit 4) Timer 0 reached its target value.
        \item[INT\_PS2\_MOUSE] (bit 3) PS/2 mouse data.
        \item[INT\_PS2\_KBD] (bit 2) PS/2 keyboard data.
        \item[INT\_VKY\_SOL] (bit 1) VICKY start of line: the display reached the line set in the line interrupt registers.
        \item[INT\_VKY\_SOF] (bit 0) VICKY start of frame.
    \end{description}
    \item[INT\_PENDING\_1] (\verb+$FE21+, read/write, reset \verb+$00+) Group 1 pending flags. Bits 3–1 have no source.
    \begin{description}
        \item[INT\_SDC\_INS] (bit 7) An SD card was inserted.
        \item[INT\_VIA1] (bit 6) VIA 1, which scans the K2's built-in keyboard. Not connected on the Jr2.
        \item[INT\_VIA0] (bit 5) VIA 0 (the joystick ports).
        \item[INT\_RTC] (bit 4) Real-time clock.
        \item[INT\_UART] (bit 0) The serial port (UART).
    \end{description}
    \item[INT\_PENDING\_2] (\verb+$FE22+, read/write, reset \verb+$00+) Group 2 pending flags: the IEC bus lines and the interrupt pins of the network, Wi-Fi and HDMI parts. Bit 7 has no source. The IEC bits can also raise an NMI (see the text).
    \begin{description}
        \item[INT\_HDMI\_PIN] (bit 6) HDMI interface interrupt pin.
        \item[INT\_WIFI\_PIN] (bit 5) Wi-Fi module interrupt pin.
        \item[INT\_NET\_PIN] (bit 4) Network module interrupt pin.
        \item[IEC\_SREQ\_i] (bit 3) IEC SREQ line.
        \item[IEC\_ATN\_i] (bit 2) IEC ATN line.
        \item[IEC\_CLK\_i] (bit 1) IEC clock line.
        \item[IEC\_DATA\_i] (bit 0) IEC data line.
    \end{description}
    \item[INT\_PENDING\_3] (\verb+$FE23+, read/write, reset \verb+$00+) Group 3 pending flags: receive and transmit FIFOs. Bit 4 (INT\_MIDI\_VS\_RX in NitrOS-9) and bits 7–6 have no source on either board.
    \begin{description}
        \item[INT\_WIZFI\_TX] (bit 5) The Wi-Fi transmit FIFO became empty.
        \item[INT\_WIZNET] (bit 3) The WizNet FIFO has data. K2 only.
        \item[INT\_OPT\_KBD] (bit 2) The optical keyboard FIFO has data. K2 only.
        \item[INT\_MIDI\_RX] (bit 1) The MIDI receive FIFO has data.
        \item[INT\_WIZFI\_RX] (bit 0) The Wi-Fi receive FIFO has data.
    \end{description}
    \item[INT\_POLARITY\_0–3] (\verb+$FE24+--\verb+$FE27+, read/write, reset \verb+$00+) Polarity for groups 0–3 (\$FE24–\$FE27), one bit per source in the same positions as the pending flags. With EDGE, chooses the condition that sets the pending flag.
    \item[INT\_EDGE\_0–3] (\verb+$FE28+--\verb+$FE2B+, read/write, reset \verb+$FF+) Edge select for groups 0–3 (\$FE28–\$FE2B). 1 = the flag is set by a change of the input; 0 = it is set while the input is at a level.
    \item[INT\_MASK\_0–3] (\verb+$FE2C+--\verb+$FE2F+, read/write, reset \verb+$FF+) Masks for groups 0–3 (\$FE2C–\$FE2F). 1 = the source can set its pending flag but doesn't interrupt the CPU; 0 = a pending flag interrupts the CPU.
\end{description}
